\chapter{Visualization and Exploratory Analysis}

\WeekHeader{7}{Visualization and Exploratory Analysis}{Use figures to reveal patterns, gaps, and hypotheses without overclaiming.}{At least three publication-style exploratory figures with captions.}

\section{Why This Matters}
Figures are where students often discover whether the research question is working. A good exploratory figure makes patterns visible, but also shows uncertainty, missing data, and source coverage. Week 7 trains students to make figures that can survive dsDNA Core review.

\section{Figure Types}
\begin{tabularx}{\textwidth}{p{0.25\textwidth}Y}
\toprule
\textbf{Figure} & \textbf{Use}\\
\midrule
Source coverage heatmap & Shows which databases contributed evidence for each chemical.\\
Trait matrix heatmap & Groups chemicals by shared discrete traits.\\
Property scatterplot & Compares measured or computed properties such as exact mass, molecular weight, XLogP, or TPSA when available.\\
Bar chart & Summarizes trait counts, hazard code counts, source hits, or pathway roles.\\
Network sketch & Shows chemical-trait or chemical-pathway relationships for a small set.\\
\bottomrule
\end{tabularx}

\section{Guided Lab}
Use the script in \path{training/scripts/07_visualization.R}. It creates example figures from a saved \texttt{categorate()} result.

\begin{rcode}
source("training/scripts/07_visualization.R")

plots <- build_training_figures(result, output_dir = "figures")
names(plots)
\end{rcode}

\section{Caption Standard}
Each figure caption should include:
\begin{itemize}
  \item What the figure shows.
  \item Which table or matrix was used.
  \item What counts as evidence.
  \item What limitation matters most.
  \item One interpretation that can be defended.
\end{itemize}

\section{Independent Extension}
Students produce a figure set that answers one of these:
\begin{itemize}
  \item Which chemicals share the most traits?
  \item Which source contributed the most useful evidence?
  \item Which chemical classes or traits separate sample groups?
  \item Which compounds need additional confirmation before interpretation?
  \item Which chemical traits connect most directly to the project question?
\end{itemize}

\section{Deliverables}
\begin{tabularx}{\textwidth}{p{0.25\textwidth}YY}
\DeliverableTableHeader
Figure set & At least three figures, each with a caption and source table. & PNG or PDF\\
Figure script & Recreates all figures from processed data. & R script\\
Exploratory interpretation & Names patterns, caveats, and next analysis decisions. & Report draft\\
\end{tabularx}

\section{Quality Checklist}
\begin{checklistbox}{Before Week 8}
\begin{itemize}
  \item Every figure can be regenerated from a script.
  \item Axes, labels, colors, and legends are readable.
  \item Captions do not claim more than the data support.
  \item Missing source coverage is visible or discussed.
\end{itemize}
\end{checklistbox}
